\newpage
\subsubsection{Описание и кодирование выходных воздействий}
\label{sec:outputs}

Выходной алфавит $Y = \{z_0,z_1,\ldots,z_{10}\}$ содержит команды, которые автомат подает при обработке входных событий. Выход $z_0$ означает работу по режиму, заданному состоянием. Например, при переходе в $S_0$ индикатор показывает текущее время, а в $S_{10}$ отображается время на момент останова.

Выходы кодируются четырьмя разрядами $y_3y_2y_1y_0$ по номеру символа (таблица~\ref{tab:outputs}).

\begin{table}[H]
	\centering
	\caption{Кодирование выходных воздействий}
	\label{tab:outputs}
	\small\setstretch{1.12}
	\renewcommand{\arraystretch}{1.10}
	\begin{tabularx}{\textwidth}{|c|>{\raggedright\arraybackslash}X|c|}
		\hline
		Выход & Действие & Код \\
		\hline
		$z_0$ & Работа по режиму, заданному состоянием & $0000$ \\
		$z_1$ & Сбросить значение секунд текущего времени до $00$ & $0001$ \\
		$z_2$ & Увеличить значение минут на единицу & $0010$ \\
		$z_3$ & Увеличить значение часов на единицу & $0011$ \\
		$z_4$ & Разрешить отсчет секундомера с сохраненного значения & $0100$ \\
		$z_5$ & Остановить секундомер и сохранить показание & $0101$ \\
		$z_6$ & Сбросить показание остановленного секундомера & $0110$ \\
		$z_7$ & Переключить признак включения будильника & $0111$ \\
		$z_8$ & Увеличить значение минут будильника на единицу & $1000$ \\
		$z_9$ & Увеличить значение часов будильника на единицу & $1001$ \\
		$z_{10}$ & Выбрать следующую мелодию будильника & $1010$ \\
		\hline
	\end{tabularx}
\end{table}

Команда $z_4$ устанавливает признак хода секундомера, а $z_5$ снимает его. Команда $z_6$ обнуляет показание остановленного секундомера. В $S_{10}$ ход счетчика текущего времени приостановлен. Четвертьчасовой звук разрешен в течение пребывания автомата в $S_{12}$, поэтому обе дуги сигнализации имеют выход $z_0$.
